\chapter{A LINGUAGEM MINIPAR 2026.2 E GRAMÁTICA BNF}

A linguagem MiniPar 2026.2 é uma linguagem imperativa, concorrente, distribuída e fortemente tipada em seu modelo conceitual, integrando elementos sintáticos inspirados em Python, C++ e Go. Sua gramática formal suporta abstrações para Orientação a Objetos, controle de fluxo abrangente, manipulação de coleções dinâmicas e estruturas para comunicação e paralelismo.

\section{Características Fundamentais e Novos Paradigmas}

O compilador dá suporte formal ao paradigma de real paralelismo com execução em múltiplos núcleos e nós de rede:

\begin{enumerate}
    \item Blocos \texttt{seq} e \texttt{par}:
    O comando \texttt{seq \{ ... \}} garante a ordem sequencial de execução dos enunciados internos. O bloco \texttt{par \{ ... \}} dispara os comandos contidos como threads concorrentes no computador local, aguardando que todas concluam antes de prosseguir.
    
    \item Primitivas \texttt{fork} e \texttt{join} (Paralelismo Real Distribuído):
    A primitiva \texttt{fork} realiza a divisão de uma tarefa em subtarefas independentes alocadas explicitamente a recursos de processamento:
    \begin{itemize}
        \item \texttt{fork core(id) comando;} ou \texttt{fork core(id) \{ ... \}}: realiza a alocação determinística da subtarefa a um núcleo de CPU do processador local, isolando sua execução via afinidade de hardware (\textit{Core Affinity}).
        \item \texttt{fork host(ip, porta) comando;}: despacha e agenda a tarefa para execução em um computador remoto via rede.
    \end{itemize}
    A primitiva complementar \texttt{join} implementa a barreira de sincronização e coleta de dados:
    \begin{itemize}
        \item \texttt{join;} : aguarda o término de todas as tarefas disparadas pelo pool local.
        \item \texttt{join(t1, t2, ...);} : bloqueia a execução até que as tarefas indicadas concluam e disponibilizem seus resultados.
    \end{itemize}

    \item Comunicação por Mensagens via Canais (\texttt{c\_channel}):
    A troca de dados entre processos independentes ocorre por passagem de mensagens utilizando canais sobre sockets TCP/IP:
    \begin{itemize}
        \item Declaração: \texttt{var canal = c\_channel("127.0.0.1", 5000);}
        \item Envio: \texttt{send(canal, mensagem);}
        \item Recepção: \texttt{var dado = receive(canal);}
    \end{itemize}

    \item Orientação a Objetos Plena:
    Suporte completo a classes (\texttt{class}), herança simples (\texttt{extends}), instanciação explícita (\texttt{new}), referência à própria instância (\texttt{this}), construtores parametrizados e métodos com passagem de parâmetros.
\end{enumerate}

\section{Gramática Formal em Notação EBNF}

A sintaxe formal completa do MiniPar 2026.2 é descrita na notação BNF Estendida (EBNF), na qual \texttt{\{ x \}} denota zero ou mais repetições, \texttt{[ x ]} denota elemento opcional e \texttt{"token"} representa símbolo terminal.

\subsection{Estrutura Global do Programa}
\begin{lstlisting}[language=bash, caption={Produções globais e declarações}]
<programa> ::= { <declaracao> | <comando> }

<declaracao> ::= <declaracao_classe>
               | <declaracao_funcao>
               | <declaracao_variavel>

<declaracao_classe> ::= "class" <IDENTIFIER> [ "extends" <IDENTIFIER> ] "{" 
                          { <membro_classe> } 
                        "}"

<membro_classe> ::= <declaracao_variavel>
                  | <declaracao_funcao>

<declaracao_funcao> ::= "func" <IDENTIFIER> "(" [ <lista_parametros> ] ")" <bloco_comandos>

<lista_parametros> ::= <IDENTIFIER> { "," <IDENTIFIER> }

<declaracao_variavel> ::= ( "var" | <tipo> ) <IDENTIFIER> [ "=" <expressao> ] ";"
\end{lstlisting}

\subsection{Comandos e Estruturas de Controle}
\begin{lstlisting}[language=bash, caption={Produções de comandos e controle de fluxo}]
<comando> ::= <bloco_comandos>
            | <bloco_seq>
            | <bloco_par>
            | <comando_fork>
            | <comando_join>
            | <comando_if>
            | <comando_while>
            | <comando_do_while>
            | <comando_for>
            | <comando_print>
            | <comando_send>
            | <comando_return>
            | <comando_break>
            | <comando_continue>
            | <atribuicao> ";"
            | <chamada_metodo> ";"
            | <chamada_funcao> ";"

<bloco_comandos> ::= "{" { <comando> } "}"
<bloco_seq>      ::= "seq" "{" { <comando> } "}"
<bloco_par>      ::= "par" "{" { <comando> } "}"

<comando_fork>   ::= "fork" <destino> ( <bloco_comandos> | <expressao> ";" )
<destino>        ::= "core" "(" <expressao> ")" 
                   | "host" "(" <expressao> [ "," <expressao> ] ")"

<comando_join>   ::= "join" [ "(" [ <lista_expressoes> ] ")" ] ";"

<comando_if>     ::= "if" "(" <expressao> ")" <comando> [ "else" <comando> ]

<comando_while>  ::= "while" "(" <expressao> ")" <comando>

<comando_do_while> ::= "do" <comando> "while" "(" <expressao> ")" ";"

<comando_for>    ::= "for" "(" ( <for_tradicional> | <for_in> ) ")" <comando>
<for_tradicional> ::= [ <declaracao_variavel> | <atribuicao> ] ";" 
                      [ <expressao> ] ";" 
                      [ <atribuicao> | <expressao> ]
<for_in>          ::= "var" <IDENTIFIER> "in" <expressao>

<comando_print>  ::= "print" "(" [ <lista_expressoes> ] ")" ";"
<comando_send>   ::= "send" "(" <expressao> "," <expressao> ")" ";"
<comando_return> ::= "return" [ <expressao> ] ";"
<comando_break>  ::= "break" ";"
<comando_continue> ::= "continue" ";"
\end{lstlisting}

\subsection{Expressões e Precedência de Operadores}
\begin{lstlisting}[language=bash, caption={Hierarquia sintática de expressões}]
<expressao> ::= <expressao_logica_ou>

<expressao_logica_ou> ::= <expressao_logica_e> { "||" <expressao_logica_e> }

<expressao_logica_e>  ::= <expressao_igualdade> { "&&" <expressao_igualdade> }

<expressao_igualdade> ::= <expressao_relacional> { ( "==" | "!=" ) <expressao_relacional> }

<expressao_relacional> ::= <expressao_aditiva> { ( "<" | "<=" | ">" | ">=" ) <expressao_aditiva> }

<expressao_aditiva>   ::= <expressao_multiplicativa> { ( "+" | "-" ) <expressao_multiplicativa> }

<expressao_multiplicativa> ::= <expressao_unaria> { ( "*" | "/" | "%" ) <expressao_unaria> }

<expressao_unaria>    ::= ( "-" | "!" ) <expressao_unaria> | <expressao_posfixa>

<expressao_posfixa>   ::= <expressao_primaria> { 
                            "." <IDENTIFIER> [ "(" [ <lista_expressoes> ] ")" ]
                          | "[" <expressao> "]" 
                        }

<expressao_primaria>  ::= <NUMBER>
                        | <STRING_LITERAL>
                        | <BOOL_LITERAL>
                        | <IDENTIFIER>
                        | "this"
                        | "new" <IDENTIFIER> "(" [ <lista_expressoes> ] ")"
                        | "c_channel" "(" <expressao> "," <expressao> ")"
                        | "receive" "(" <expressao> ")"
                        | "input" "(" [ <expressao> ] ")"
                        | "matrix" "(" <expressao> "," <expressao> "," <expressao> ")"
                        | "[" [ <lista_expressoes> ] "]"
                        | "(" <expressao> ")"
                        | <comando_fork>
                        | <comando_join>
\end{lstlisting}
